\documentclass[11pt,a4paper]{article}
\usepackage[utf8]{inputenc}
\usepackage[T1]{fontenc}
\usepackage[ngerman]{babel}
\usepackage{amsmath,amssymb}
\usepackage[margin=2.3cm]{geometry}
\usepackage{enumitem}
\usepackage{array}
\usepackage{booktabs}
\usepackage{xcolor}
\usepackage{titlesec}

\titleformat{\section}{\large\bfseries\color{blue!55!black}}{Aufgabe \thesection}{0.6em}{}
\setlist{nosep}

\newcommand{\pt}[1]{\hfill\textnormal{\small\color{gray}(#1\,P)}}
\newcommand{\korr}[1]{\par\smallskip{\color{red}\textbf{Korrektur:} #1}}

\begin{document}

\begin{center}
  {\Large\bfseries INF4:1 -- Übungsblatt Kapitel 02}\\[2pt]
  {\normalsize Effizienz \& Landau-Symbole \quad(passend zu Praktikum 1)}\\[2pt]
  {\small\itshape Erst selbst lösen -- Lösungen besprechen wir danach gemeinsam.}
\end{center}
\vspace{0.6em}

%--------------------------------------------------
\section{O / \texorpdfstring{$\Omega$}{Omega} / \texorpdfstring{$\Theta$}{Theta} ankreuzen \pt{6}}
Wahr oder falsch? \textbf{Minuspunkte für Falsches} -- im Zweifel weglassen (wie in der Klausur).
\smallskip

\renewcommand{\arraystretch}{1.5}
\begin{tabular}{@{}c >{$}l<{$} c@{}}
\toprule
 & \text{Aussage} & wahr / falsch \\
\midrule
a) & n \in O(n^2) & true \\
b) & n^2 \in O(n) & false \\
c) & n^2 \in \Omega(n) & true \\
d) & 3n^2 + 5n - 7 \in \Theta(n^2) & true \\
e) & 2^n \in O(n^2) & {\color{red}true $\rightarrow$ \textbf{false}} \\
f) & \log_2 n \in \Omega(n) & false \\
\bottomrule
\end{tabular}

\smallskip
\textit{Achtung Richtung: } $O$ = obere Schranke (wächst \emph{höchstens} so schnell),
$\Omega$ = untere Schranke (wächst \emph{mindestens} so schnell).

\korr{a)--d), f) richtig. \textbf{e) falsch:} $2^n \notin O(n^2)$ -- exponentiell schlägt
jedes Polynom. Wachstumshierarchie: $1 < \log n < n < n\log n < n^2 < n^3 < 2^n < n!$.}

%--------------------------------------------------
\section{Komplexität bestimmen \pt{4}}
Gib die \textbf{$\Theta$-Klasse} in Abhängigkeit von $n$ an:
\begin{verbatim}
    s = 0
    for i = 1 to n:
        for j = i to n:
            s = s + 1
\end{verbatim}
Wie oft wird \texttt{s = s + 1} insgesamt ausgeführt (geschlossener Ausdruck)?
Und \emph{warum} läuft die innere Schleife nicht $n$-mal pro Durchlauf?
\\
Antwort: $(n(n+1))/2$-Durchläufe. Die innere Schleife laeft haeufiger als n mal durch, da sie pro Durchlauf der aeusseren Schleife mehrmals (mind. 1x) durchlaufen wird.

\korr{Geschlossener Ausdruck $\tfrac{n(n+1)}{2}$ \emph{richtig}; Klasse fehlt noch:
$\tfrac{n^2+n}{2}\to\Theta(n^2)$. Begründung umgekehrt: die innere Schleife läuft
\emph{seltener}, nicht häufiger, als $n$-mal -- sie macht $n-i+1$ Durchläufe, also
\emph{abnehmend} von $n$ (bei $i=1$) auf $1$ (bei $i=n$), im Schnitt $\approx n/2$.}

%--------------------------------------------------
\section{$c$ und $n_0$ finden \pt{3}}
Zeige formal $f(n) \in O\big(g(n)\big)$, indem du \textbf{konkrete} Werte für $c$ und $n_0$ angibst,
sodass $0 \le f(n) \le c\cdot g(n)$ für alle $n \ge n_0$ gilt:
\[
  f(n) = 4n + 10, \qquad g(n) = n .
\]
\newline
Rechnung auf Papier.
$4n=10 <= 4n+10n = 14n$\\
$c=14, n_0=1$
$f(n) \in O\big(g(n)\big)$ für $c=14, n>=1$

\korr{Richtig! $4n+10 \le 4n+10n = 14n$ für $n\ge 1$, also $c=14,\ n_0=1$. \;(kleiner
Tippfehler oben: gemeint ist $4n+10$, nicht $4n=10$.)}

%--------------------------------------------------
\section{Algorithmus entwerfen \pt{8}}
Gegeben ein Array mit $n$ Zahlen. Beschreibe in \textbf{Pseudocode} einen Algorithmus, der prüft,
ob das Array ein \textbf{Duplikat} enthält -- mit Laufzeit \textbf{$\Theta(n\log n)$}.
Welcher Trick bringt dich von der naiven $\Theta(n^2)$-Lösung herunter?
////
Geht man hier davon aus dass das Array schon sortiert ist? Kommt erst im naechstem Kapitel oder?

\korr{Nein -- du \emph{sortierst selbst}. Trick: (1)~Array mit einem bekannten
$\Theta(n\log n)$-Verfahren (z.\,B.\ Merge Sort) sortieren -- \emph{wie} das geht kommt
erst in Kap.~03/04, du darfst es als Black Box benutzen. (2)~\emph{Ein} linearer Durchlauf
($\Theta(n)$): vergleiche je zwei Nachbarn $a[i]$ und $a[i+1]$ -- Duplikate stehen nach dem
Sortieren direkt nebeneinander. Gesamt $\Theta(n\log n)+\Theta(n)=\Theta(n\log n)$. Das
schlägt die naive Variante (jedes Paar prüfen, $\Theta(n^2)$).}
%--------------------------------------------------
\section{$\Theta$ oder nur $O$? \pt{6}}
Entscheide jeweils, ob $f(n) \in \Theta\big(g(n)\big)$ \emph{oder nur} $f(n) \in O\big(g(n)\big)$ gilt,
und begründe kurz:
\begin{enumerate}[label=\alph*)]
  \item $f(n) = 5n^2 + 3n,\qquad g(n) = n^2$
  \item $f(n) = n,\qquad g(n) = n^2$
  \item $f(n) = n \cdot \log_2 n,\qquad g(n) = n^2$
\end{enumerate}
\begin{enumerate}[label=\alph*)]
  \item $f(n) \in \Theta\big(g(n)\big)$ da $5n^2 >= n^2$
  \item Nur $f(n) \in O\big(g(n)\big)$ da mit $n>=cn^2$ kein ganzzahliger Wert fuer C gefunden werden kann.
  \item Nur $f(n) \in O\big(g(n)\big)$ da kein eindeutiger Wert fuer C gefunden werden kann.
\end{enumerate}

\korr{Alle drei Schlüsse richtig (a)~$\Theta$, b)~nur $O$, c)~nur $O$). Begründung schärfen:
$c$ muss \textbf{nicht} ganzzahlig/eindeutig sein. „nur $O$`` scheitert an der \textbf{unteren}
Schranke $\Omega$, weil $f$ \emph{strikt langsamer} wächst, d.h.\ $f(n)/g(n)\to 0$:
b)~$\tfrac{n}{n^2}=\tfrac1n\to0$;\; c)~$\tfrac{n\log n}{n^2}=\tfrac{\log n}{n}\to0$.}

\vfill
\hrule
\smallskip
{\footnotesize \textit{Bonus (Kap.~04-Vorgriff):} Warum ist Merge Sort $\Theta(n\log n)$, und warum
gibt es dort keinen Unterschied zwischen Best- und Worst-Case? Ist Merge Sort in-place / stabil?}

\end{document}
